\section{Truth, satisfaction, validity, and models}
\label{sec:levels-04}

The notions of structure, variable assignment, sentence, and first-order
formula used in this section receive their exact inductive definitions in
\cref{ch:fol-syntax-semantics} and \cref{ch:fol-semantics}.  The
definitions below are stated here schematically, to fix the registered
senses of the semantic vocabulary; nothing before those chapters depends on
more than the registered senses.

\begin{definition}[Satisfaction]
\emph{Satisfaction} is the relation defined by the semantic clauses of the
language.  In first-order logic it has the form
\(\mathcal M,s\entails\varphi\), where \(\mathcal M\) is a structure,
\(s\) is a variable assignment, and \(\varphi\) is a formula.
\end{definition}

\begin{definition}[Truth in a structure]
A sentence \(\sigma\) is \emph{true in} a structure \(\mathcal M\) iff
\(\mathcal M,s\entails\sigma\) for every variable assignment \(s\).
Because a sentence has no free variable occurrence, the choice of \(s\) will
later be proved irrelevant.  We then write \(\mathcal M\entails\sigma\).
\end{definition}

\begin{definition}[Model]
Let \(T\) be a set of sentences.  A structure \(\mathcal M\) is a
\emph{model of} \(T\) iff \(\mathcal M\entails\sigma\) for every
\(\sigma\in T\).  The class of models of \(T\) is
\[
  \Mod(T)\defeq\set{\mathcal M: \mathcal M\entails T},
\]
where \(\mathcal M\entails T\) abbreviates: \(\mathcal M\entails\sigma\)
for every \(\sigma\in T\).
\end{definition}

Thus a structure is not called a model without an accompanying theory or
sentence.  The phrase ``the model \(\mathcal M\)'' is permitted only after the
relevant theory has been fixed.

